\documentclass[10pt]{article}
\input{common}
\title{RevIN Experiment 3: Centering and Transform Variants}
\author{}
\date{}
\begin{document}\maketitle
\section{Theoretical overview and goal}
Mean centering assumes the look-back average is the appropriate location anchor
for the future.  Last-value centering instead predicts changes around the most
recent observation, which may better suit local trends.  The inverse hyperbolic
sine transform
\[
 z=\operatorname{asinh}(x)
\]
is approximately linear near zero and logarithmic for large magnitudes, giving
a reversible compression of large standardized magnitudes after instance
normalization.

\goal{Test two lightweight changes to RevIN's inductive bias: a more local
location anchor and a nonlinear magnitude compression as appendix ablations,
without confounding either transformation with the training-loss choice.}

\section{Implementation details}
The variants use non-affine instance normalization and are each trained with
both MSE and normalized MSE:
\begin{center}
\begin{tabular}{llll}
\toprule
Variant & Center & Transform & Losses \\
\midrule
\code{instance} & Look-back mean & Identity & MSE, nMSE \\
\code{revin\_last} & Last observation & Identity & MSE, nMSE \\
\code{revin\_arcsinh} & Look-back mean & asinh & MSE, nMSE \\
\bottomrule
\end{tabular}
\end{center}
The concrete method names append \code{\_mse} or \code{\_nmse}; for example,
\code{revin\_last\_mse} and \code{revin\_last\_nmse}. The historical
\code{revin\_last} and \code{revin\_arcsinh} labels denote non-affine instance
normalization despite retaining the RevIN prefix.
Every transformation is inverted before data-space evaluation. These appendix
variants run in the full and ultra profiles, over the same nine datasets, five
settings, PatchTST (plus DLinear in ultra), three seeds, optimizer, batch size, 10,000-step
budget, and validation frequency. Configuration mapping
is defined by \code{method\_args} in
\code{src/slurm/run\_revin\_experiment.sh}; layer behavior is in
\code{src/utils/models.py}.
The root front runs the separate training and table stages in order and resumes
from signature-matched completion markers.

Tables are separate by model and test split.  They report ordinary MSE as mean
$\pm$ sample standard deviation with explicit $\times10^m$ row scaling. Each
variant is compared with \code{instance\_mse} under MSE training and
\code{instance\_nmse} under nMSE training.

\section{Interpretation}
Last centering tests sensitivity to local level evolution; asinh tests
sensitivity to tails and multiplicative scale.  A gain is evidence for that
specific preprocessing bias, not a general solution to conditional shift.
Results should be examined across datasets because either transform may remove
useful absolute-level information.
\end{document}
